Uma \emph{relação}~$R$ sobre um conjunto~$A$ é um modo de relacionar elementos
de~$A$. Escrevemos $Rxy$ se a relação \emph{vale} entre $x$
e~$y$. Formalmente, podemos considerar $R$ como o conjunto de pares
$\tuple{x,y} \in A^2$ tais que~$Rxy$. Ser menor que, maior que,
igual a, dividir exatamente, ter o mesmo comprimento que, ser um subconjunto de, e ter o
mesmo tamanho que são todos exemplos importantes de relações (sobre
conjuntos de números, cadeias, ou de conjuntos). \emph{Grafos} são um modo geral
de representar relações visualmente. Mas um grafo também pode ser visto como uma
relação binária (a relação \emph{aresta}) junto com o
conjunto subjacente de \emph{vértices}.

Algumas relações compartilham certas características que as tornam especialmente
interessantes ou úteis. Uma relação~$R$ é \emph{reflexiva} se
tudo está $R$-relacionado consigo mesmo; \emph{simétrica}, se, com $Rxy$,
também $Ryx$ vale para quaisquer $x$ e $y$; e \emph{transitiva} se $Rxy$
e $Ryz$ garantem~$Rxz$. As relações que têm todas essas três
propriedades são \emph{relações de equivalência}. Uma relação é
\emph{anti-simétrica} se $Rxy$ e $Ryx$
garantem~$x=y$. \emph{Ordens parciais} são aquelas relações que são
reflexivas, anti-simétricas e transitivas. Uma \emph{ordem linear} é
qualquer ordem parcial que satisfaz que, para quaisquer $x$ e~$y$, ou
$Rxy$ ou $x=y$ ou $Ryx$. (Em geral, uma relação com essa propriedade é
\emph{conexa}).

Como as relações são conjuntos (de pares), elas podem ser operadas como conjuntos
(por exemplo, podemos formar a união e a interseção de relações). Também podemos
encadeá-las (\emph{produto relativo}~$R \mid S$). Se
formarmos o produto relativo de~$R$ consigo mesmo um número arbitrário de vezes,
obtemos o \emph{fecho transitivo}~$R^+$ de~$R$.
